\documentclass[12pt]{exam}
\usepackage[T1]{fontenc}
\usepackage[utf8]{inputenc}
\usepackage[lithuanian]{babel}
\usepackage{lmodern}
\usepackage{graphicx}
\usepackage[a4paper, left=1cm, right=1cm, top=2.5cm, bottom=1cm]{geometry}
\usepackage{amsmath, amssymb}
\usepackage{tasks}

% Antraštė
\pagestyle{headandfoot}
\runningheadrule
\firstpageheader{Vardas: \rule{4cm}{0.4pt}}{Kontrolinis darbas: \\Daugianariai}{Šviesos variantas}
\runningheader{Daugianariai}{}{}
\firstpagefooter{}{}{}
\runningfooter{}{}{}
\pointsdroppedatright          
\bracketedpoints               
\setlength{\rightpointsmargin}{2.5cm} 
\renewcommand{\points}{taškai} 
\renewcommand{\pointname}{Taškai:} 
\marginpointname{\ taškai}     
\renewcommand{\dottedlinefillheight}{0.7cm}
\newcommand{\atsakymas}{%
    \par\vspace{0.3cm}
    \textbf{Atsakymas:} \framebox[7cm]{\rule{0pt}{0.7cm}}
}
\renewcommand{\partshook}{%
    \setlength{\leftmargin}{1pt}%
    \setlength{\itemindent}{17pt}%
    \setlength{\labelsep}{0.5em}%
}

\begin{document}

\begin{questions}

\question[6] Nustatykite kiekvieno reiškinio tipą: racionalusis (R), trupmeninis (T), sveikasis (S) ar iracionalusis (I). Prie kiekvieno reiškinio nurodykite visas jam tinkančias raides.
\droppoints
\begin{tasks}[label-width=1.2em](3)
\task $\dfrac{3x-2}{\sqrt7}-x$
\task $5x^{-3}+2x-9$
\task $\dfrac{x^2+1}{x+2}+\sqrt3$
\task $\sqrt{7-2x}+3x-4$
\task $16^{-5}+x^2+2x-\sqrt{10}$
\task $\dfrac5{x^2}-3x+8$
\end{tasks}



\question[5] Suprastinkite reiškinį ir nurodykite, ar gautas reiškinys yra vienanaris. Jei reiškinys yra vienanaris, nurodykite jo koeficientą ir laipsnį; jei ne, paaiškinkite kodėl.

\droppoints
    $$\displaystyle\frac{(-2x^{11}y^{-3})^8\cdot(16x^{-2}y^5)^{-6}}{(32x^8y^{-9})^5}$$

\vspace{3cm}
\atsakymas

\question Suprastinkite reiškinį $(2x-1)^2(x+1)-(x-1)^3-3(x+2)(x-1)(x+3)+8x ^{2} -19
$ ir užrašę jį pavidalu $a(x+b)^2+c$ raskite:
\begin{parts}
    \part[7] didžiausią arba mažiausią jo reikšmę bei nurodykite $x$ reikšmę su kuria ta reikšmė gaunama;
   
    \droppoints
    \vspace{10cm}
    \atsakymas
    \newpage
    \part[3] apskaičiuokite $b \cdot \left(a ^{1\overbrace{0\ldots 0}^{666}2} + 21,25 : c\right)    $ .
    \droppoints
    \vspace{2cm}
    \atsakymas
\end{parts}


\question[5] Raskite parametrus $a$, $b$ ir $c$ kad kiekviena lygybė galiotų su visomis $x$ reikšmėmis.
\droppoints
$$(ax+a)(3x^2-bx+c)=-6x^3+4x^2+2x-8.$$

\vspace{5cm}
\atsakymas

\question[4] Raskite žemiau esančio reiškinio mažiausią reikšmę:
\droppoints
$$4x^2+9y^2+z^2+12x-12y-10z+40.$$

\vspace{5cm}
\atsakymas

\end{questions}

\vfill 

\begin{center}
\textbf{Vertinimo lentelė (iš viso 30 taškų)}\par\medskip
\setlength{\tabcolsep}{5pt}
\renewcommand{\arraystretch}{1.3}
\begin{tabular}{|c|*{10}{c|}}
\hline
Taškai & 0--3 & 4--6 & 7--9 & 10--12 & 13--15 & 16--18 & 19--22 & 23--25 & 26--28 & 29--30\\\hline
Pažymys & 1 & 2 & 3 & 4 & 5 & 6 & 7 & 8 & 9 & 10\\\hline
\end{tabular}
\end{center}

\end{document}
